% !TEX program = xelatex
\documentclass[10pt,letterpaper,twocolumn]{article}

\usepackage[top=0.75in, bottom=0.75in, left=0.75in, right=0.75in]{geometry}
\usepackage{fontspec}
\usepackage{unicode-math}
\setmainfont{Latin Modern Roman}[
  Extension=.otf,
  UprightFont=lmroman10-regular,
  BoldFont=lmroman10-bold,
  ItalicFont=lmroman10-italic,
  BoldItalicFont=lmroman10-bolditalic,
]
\setsansfont{Latin Modern Sans}[
  Extension=.otf,
  UprightFont=lmsans10-regular,
  BoldFont=lmsans10-bold,
  ItalicFont=lmsans10-oblique,
  BoldItalicFont=lmsans10-boldoblique,
]
\newfontfamily\acbold{ywft-processing-bold}[Path=../../system/public/type/webfonts/,Extension=.ttf]
\newfontfamily\aclight{ywft-processing-light}[Path=../../system/public/type/webfonts/,Extension=.ttf]
\setmonofont{Latin Modern Mono}[
  Extension=.otf,
  UprightFont=lmmono10-regular,
  ItalicFont=lmmono10-italic,
  Scale=0.85,
]

\usepackage{xcolor}
\usepackage{titlesec}
\usepackage{enumitem}
\usepackage{booktabs}
\usepackage{tabularx}
\usepackage{fancyhdr}
\usepackage{hyperref}
\usepackage{xurl}
\usepackage{graphicx}
\graphicspath{{figures/}{../../papers/arxiv-ac/figures/}}
\usepackage{ragged2e}
\usepackage{microtype}
\usepackage{natbib}
\usepackage[colorspec=0.92]{draftwatermark}

\definecolor{acpink}{RGB}{180,72,135}
\definecolor{acpurple}{RGB}{120,80,180}
\definecolor{acdark}{RGB}{64,56,74}
\definecolor{acgray}{RGB}{119,119,119}
\definecolor{draftcolor}{RGB}{180,72,135}

\DraftwatermarkOptions{text=WORKING DRAFT,fontsize=3cm,color=draftcolor!18,angle=45}

\hypersetup{colorlinks=true,linkcolor=acpurple,urlcolor=acpurple,citecolor=acpurple,
  pdftitle={Sucking on the Complex: Hegemon\'{i}a de plataformas, cr\'{i}tica-como-contenido y la necesidad de anti-entornos}}

\titleformat{\section}{\normalfont\bfseries\normalsize\uppercase}{\thesection.}{0.5em}{}
\titlespacing{\section}{0pt}{1.2em}{0.3em}
\titleformat{\subsection}{\normalfont\bfseries\small}{\thesubsection}{0.5em}{}
\titlespacing{\subsection}{0pt}{0.8em}{0.2em}

\pagestyle{fancy}\fancyhf{}
\renewcommand{\headrulewidth}{0pt}
\fancyhead[C]{\footnotesize\color{acpink}\textit{Borrador de trabajo --- no citar}}
\fancyfoot[C]{\footnotesize\thepage}

\newcommand{\ac}{\textsc{Aesthetic.Computer}}
\newcommand{\np}{\textsc{notepat}}

\setlist[itemize]{nosep, leftmargin=1.2em, itemsep=0.1em}
\setlength{\columnsep}{1.8em}
\setlength{\parindent}{1em}
\setlength{\parskip}{0.3em}

\begin{document}

\twocolumn[{%
\begin{center}

\includegraphics[height=4em]{pals}\par\vspace{0.5em}
{\acbold\fontsize{22pt}{26pt}\selectfont\color{acdark} Sucking on the Complex}\par
\includegraphics[width=0.5\textwidth]{figures/cover}\par
\vspace{0.35em}
\vspace{0.2em}
{\aclight\fontsize{11pt}{13pt}\selectfont\color{acpink} Hegemon\'{\i}a de plataformas, cr\'{\i}tica-como-contenido y la necesidad de anti-entornos}\par
\vspace{0.6em}
{\normalsize\href{https://prompt.ac/@jeffrey}{@jeffrey}}\par
{\small\color{acgray} Aesthetic.Computer}\par
{\small\color{acgray} ORCID: \href{https://orcid.org/0009-0007-4460-4913}{0009-0007-4460-4913}}\par
\vspace{0.3em}
{\small\color{acpurple} \url{https://aesthetic.computer}}\par
\vspace{0.6em}
\rule{\textwidth}{1.5pt}
\vspace{0.5em}
\end{center}

\begin{center}
{\small\color{acpink}\textbf{[ borrador de trabajo --- no citar ]}}
\end{center}
\vspace{0.3em}

\begin{quote}
\small\noindent\textbf{Resumen.}
La infraestructura medi\'{a}tica del mundo del arte contempor\'{a}neo est\'{a} dominada por dos empresas: Meta y ByteDance. Para la mayor\'{\i}a de los artistas en activo, una presencia en plataformas es condici\'{o}n previa para la legibilidad cultural, y la cr\'{\i}tica de las plataformas se ha convertido en uno de los formatos de contenido m\'{a}s celebrados que las plataformas alojan. Un borrador anterior de este art\'{\i}culo argumentaba que los \emph{anti-entornos}---t\'{e}rmino de McLuhan para las contra-estructuras construidas fuera de un entorno dominante---son la \'{u}nica escapatoria de esta trampa, y revisaba a artistas prominentes cuyo trabajo cr\'{\i}tico con las plataformas es absorbido por las instituciones que lo exhiben. Ese argumento permanece aqu\'{\i}, comprimido. El borrador actual invierte el peso del art\'{\i}culo: en lugar de repetir el diagn\'{o}stico de la hegemon\'{\i}a de plataformas o el fracaso estructural de la cr\'{\i}tica-como-contenido institucional, reporto desde cinco a\~{n}os de construcci\'{o}n de una alternativa---\ac{}~\citep{scudder2026ac}, un runtime, red social, instrumento y sistema operativo. La afirmaci\'{o}n que este trabajo sostiene es m\'{a}s estrecha y m\'{a}s concreta que el rechazo total. Las plataformas no son adversario ni hogar; son motores de distribuci\'{o}n. El movimiento interesante para un artista en 2026 no es rechazar el motor sino usarlo deliberadamente---entregando una carga \'{u}til de notaci\'{o}n gr\'{a}fica y ret\'{o}rica de interfaz que la plataforma no puede inspeccionar, hacia un \emph{usuariado} que la plataforma no puede contener. Whistlegraph~\citep{scudder2026whistlegraph} es el caso. \ac{} es el flujo descendente. Catalogo seis compromisos estructurales sobre los que descansa un anti-entorno para la computaci\'{o}n creativa---direccionamiento URL-primero, unidades acotadas pieza-no-publicaci\'{o}n, interfaces instrumento-no-perfil, un lenguaje de la era de la composici\'{o}n, hardware de nodo soberano y comunidad de afinidad-no-seguidores---y sostengo que el trabajo de construir un anti-entorno no es el trabajo de rechazar la plataforma sino el de preparar un usuariado al que la plataforma puede distribuir, pero no contener.
\end{quote}
\vspace{0.5em}
}]

\section{El motor y el exterior}

El encuadre est\'{a}ndar de la cr\'{\i}tica de plataformas dice: el complejo est\'{a} totalizado, el artista que lo alimenta es c\'{o}mplice, y el movimiento honesto es el rechazo. Este art\'{\i}culo fue escrito originalmente en ese registro. La versi\'{o}n actual no. Cinco a\~{n}os de construcci\'{o}n de \ac{}~\citep{scudder2026ac} junto a una d\'{e}cada de gesti\'{o}n de @whistlegraph~\citep{scudder2026whistlegraph} han ense\~{n}ado que el registro est\'{a} equivocado de una manera espec\'{\i}fica: malinterpreta para qu\'{e} sirven las plataformas, desde el lado del artista.

Las plataformas son motores. Mueven atenci\'{o}n a una velocidad y volumen que ninguna tuber\'{\i}a institucional se acerca, y lo hacen de forma suficientemente fiable como para que un artista pueda tratarlas como infraestructura, del mismo modo que trata la electricidad o una ruta postal. Un artista al que le importa construir una pr\'{a}ctica que sobreviva a la plataforma tiene una opci\'{o}n que la literatura cr\'{\i}tica casi nunca describe: montar el motor deliberadamente, usando su alcance para preparar a un p\'{u}blico en una gram\'{a}tica formal que la plataforma ni escribi\'{o} ni puede reproducir, y enrutar a ese p\'{u}blico hacia una pr\'{a}ctica y un \emph{usuariado} que viven en otra pila. Esto no es rechazo. Se parece m\'{a}s a extracci\'{o}n: usar la tuber\'{\i}a para entregar una carga que la tuber\'{\i}a no puede inspeccionar.

Whistlegraph fue una de esas cargas. La forma emergi\'{o} de Radical Digital Painting~\citep{scudder2017manifesto}, 65 conferencias-performance pre-pand\'{e}micas incluyendo la gira europea de 2018 con Goodiepal \& Pals~\citep{scudder2026goodiepal} y el 35.\textsuperscript{o} Chaos Communication Congress~\citep{scudder2018rdp35c3}, y una d\'{e}cada de pr\'{a}ctica dibujo-como-canto anterior a la cuenta. TikTok no hizo viral al whistlegraph; el whistlegraph se hizo viral por s\'{\i} mismo en TikTok porque la propiedad estructural de la forma---reproducibilidad, correspondencia marca-a-s\'{\i}laba, la partitura que ense\~{n}a c\'{o}mo tocarla~\citep{scudder2026whistlegraph}---ya estaba all\'{\i}. Los 2,6 millones de seguidores de la cuenta no son la obra. Son el p\'{u}blico preparado para lo que viene tras la fase de distribuci\'{o}n de la forma: un \emph{usuariado}---una poblaci\'{o}n organizada en torno a un sistema de computaci\'{o}n creativa que puede ejecutar, bifurcar y performar dentro, en lugar de un p\'{u}blico organizado en torno a contenido que s\'{o}lo puede scrolear.

Lo que sigue no es ni una revisi\'{o}n de la hegemon\'{\i}a de plataformas (ese diagn\'{o}stico est\'{a} bien ensayado; \S2 lo comprime) ni una nueva acusaci\'{o}n a la cr\'{\i}tica-como-contenido institucional (\S3). Es un informe de lo que cinco a\~{n}os de construcci\'{o}n de la pila descendente han producido, de los compromisos estructurales sobre los que descansa un anti-entorno para la computaci\'{o}n creativa, de c\'{o}mo el motor prepara un usuariado para ello, y de lo que cuesta mantener el c\'{\i}rculo abierto. El art\'{\i}culo que est\'{a}s leyendo es \'{e}l mismo uno de esos compromisos. Est\'{a} escrito desde dentro del c\'{\i}rculo, cita el resto de la documentaci\'{o}n del c\'{\i}rculo~\citep{scudder2026ac, scudder2026kidlisp, scudder2026notepat, scudder2026os, scudder2026pieces, scudder2026urltradition, scudder2026archaeology, scudder2026whistlegraph, scudder2026goodiepal}, y se distribuye a trav\'{e}s del \emph{platter} de investigaci\'{o}n que el c\'{\i}rculo ha producido.

\section{El complejo}

Meta (Facebook, Instagram, Threads, WhatsApp) y ByteDance (TikTok, Douyin) median conjuntamente la visibilidad cultural de casi todos los artistas en activo del mundo. Instagram por s\'{\i} solo tiene m\'{a}s de tres mil millones de usuarios activos mensuales; es la plataforma principal de descubrimiento para galer\'{\i}as, curadores, coleccionistas e instituciones. El formato de v\'{\i}deo corto de TikTok se ha convertido en el mecanismo de distribuci\'{o}n predeterminado para la performance, el proceso de estudio y la educaci\'{o}n art\'{\i}stica. Juntos constituyen lo que \citet{nieborg2018platformization} llaman la ``plataformizaci\'{o}n de la producci\'{o}n cultural'': la mercanc\'{\i}a cultural contingente, conformada no por la intenci\'{o}n art\'{\i}stica sino por la l\'{o}gica de clasificaci\'{o}n de la plataforma. Para los artistas emergentes y de media carrera, las m\'{e}tricas de plataforma---seguidores, tasa de interacci\'{o}n, frecuencia de publicaciones---se han convertido en indicadores sustitutos de relevancia cultural. El precio de la visibilidad es datos: cada interacci\'{o}n alimenta la m\'{a}quina publicitaria que financia la plataforma~\citep{zuboff2019surveillance, srnicek2017platform}.

El algoritmo no muestra pasivamente el arte. Da forma activamente a qu\'{e} arte se produce. Cuando un artista publica una pintura y recibe alta interacci\'{o}n, el algoritmo la muestra a m\'{a}s personas; el artista, al ver la respuesta, hace m\'{a}s pinturas similares. Esto no es censura---es empuje conductual, el mecanismo central que \citet{zuboff2019surveillance} identifica como el n\'{u}cleo del capitalismo de vigilancia: la predicci\'{o}n y modificaci\'{o}n del comportamiento humano a gran escala, con fines de lucro. Y en Instagram, una pintura de Richter ocupa el mismo cuadrado de 1080 p\'{\i}xeles que una fotograf\'{\i}a del almuerzo de alguien. La plataforma no distingue entre el trabajo de d\'{e}cadas y el trabajo de segundos. Esto no es un fallo de curadur\'{\i}a; es el dise\~{n}o de la plataforma~\citep{chun2016updating}. La interacci\'{o}n es la \'{u}nica m\'{e}trica. \citet{lovink2011networks} llama a la estructura resultante la ``red sin causa'': un sistema que conecta todo y contextualiza nada. \citet{noble2018algorithms} y \citet{benjamin2019race} muestran que el aparato de clasificaci\'{o}n bajo el aplanamiento no es neutral; codifica y amplifica jerarqu\'{\i}as existentes. \citet{terranova2000free} identific\'{o} la consecuencia hace veinticinco a\~{n}os como ``trabajo gratis'': producci\'{o}n cultural no remunerada que genera valor para las plataformas. Los artistas son los trabajadores gratis m\'{a}s dedicados de la econom\'{\i}a de la atenci\'{o}n.

El marco m\'{a}s profundo para esto es el de Attali. En \emph{Noise}~\citep{attali1985noise} argumenta que la orquesta de sala de conciertos es un instrumento del poder del Estado: organiza los cuerpos en posiciones jer\'{a}rquicas bajo la autoridad de un director, exige obediencia s\'{\i}ncrona a una partitura escrita y produce un output unificado a partir de trabajo individual controlado. Esta es la era de la \emph{representaci\'{o}n} en la m\'{u}sica---la m\'{u}sica como espect\'{a}culo, como prueba de que el poder puede organizar el ruido en armon\'{\i}a. La plataforma tecnol\'{o}gica es la orquesta del siglo XXI. El fundador es el director. El algoritmo es la partitura. Usuarios, creadores y artistas son los m\'{u}sicos---cada uno tocando su parte, cada uno creyendo que se expresa, cada uno produciendo valor que se acumula en la instituci\'{o}n. La ret\'{o}rica de la plataforma de ``dar voz a todos'' es lenguaje de era-de-representaci\'{o}n vestido con ropa de era-de-composici\'{o}n. Todos tienen voz, pero el algoritmo decide qui\'{e}n es escuchado. Attali glorificaba lo que deb\'{\i}a venir despu\'{e}s de la representaci\'{o}n---la \emph{composici\'{o}n}, un modo donde se disuelve la distinci\'{o}n entre int\'{e}rprete y audiencia, donde la m\'{u}sica se hace por el gozo de hacerla. La plataforma no es composici\'{o}n. Es representaci\'{o}n a escala planetaria, con la extracci\'{o}n adicional de datos conductuales que ni siquiera la orquesta de la corte alcanz\'{o}. La startup tecnol\'{o}gica est\'{a} estructuralmente m\'{a}s cerca de la orquesta de la corte de los Habsburgo que de la canci\'{o}n folcl\'{o}rica. La diferencia es que la corte pagaba a sus m\'{u}sicos.

El diagn\'{o}stico no es original y no necesita serlo. Lo que sigue s\'{\i}.

\section{Cr\'{\i}tica-como-contenido}

La respuesta art\'{\i}stica dominante a la hegemon\'{\i}a de plataformas ha sido hacer obras \emph{sobre} ella. Esta obra se exhibe en las instituciones culturales m\'{a}s prestigiosas del mundo. Se celebra, se colecciona y se discute. Las condiciones estructurales que describe permanecen inalteradas.

El modo adopta varias formas. Paglen fotograf\'{\i}a instalaciones militares clasificadas y documenta la infraestructura f\'{\i}sica de la vigilancia, exhibiendo en el Whitney, MoMA, el Smithsonian y la Bienal de Venecia; la obra anuncia el alcance del complejo sin enrutar a su p\'{u}blico hacia nada m\'{a}s. \citet{steyerl2017duty} describe un mundo de ``arte libre de impuestos'' circulando por zonas de libre comercio, y sus videoinstalaciones en el Serpentine, la Bienal de Venecia y el Park Avenue Armory son an\'{a}lisis de la misma circulaci\'{o}n---compartidos en Instagram, discutidos en TikTok, la cr\'{\i}tica distribuida por la m\'{a}quina que describe (\citet{steyerl2009poor} argumentaba que la ``imagen pobre'' gana poder pol\'{\i}tico por su circulaci\'{o}n en baja resoluci\'{o}n; en el feed, cada imagen genera los mismos datos conductuales). \citet{blas2014informatic} desarroll\'{o} la ``Facial Weaponization Suite'' (2011--2014), m\'{a}scaras colectivas que derrotan al reconocimiento facial, mostradas en la Whitechapel y Tate Modern---pero la m\'{a}scara es un objeto art\'{\i}stico en una vitrina, no una herramienta que se pueda llevar. El \emph{Cyberfeminism Index} de Mindy Seu~\citep{seu2023cyberfeminism} cataloga m\'{a}s de 700 recursos que rastrean la resistencia ciberfeminista, mucho de ello distribuido a trav\'{e}s de Instagram Stories: un archivo feminista de resistencia digital hecho legible para el mundo del arte a trav\'{e}s de la infraestructura de vigilancia de Meta.

\citet{ulman2014excellences} es el caso can\'{o}nico. Su performance de cinco meses en Instagram en 2014---una manipulaci\'{o}n guionizada de la l\'{o}gica de aspiraci\'{o}n y auto-exhibici\'{o}n de la plataforma, luego exhibida en Tate Modern en 2016---tuvo \'{e}xito precisamente porque era indistinguible del contenido. La plataforma no sab\'{\i}a que estaba siendo usada como material. Los datos conductuales que gener\'{o} eran id\'{e}nticos a los de publicaciones de estilo de vida aut\'{e}nticas. Este es el punto estructural, no una acusaci\'{o}n contra ning\'{u}n artista individual: el algoritmo no distingue cr\'{\i}tica de respaldo. La fotograf\'{\i}a de Paglen de una instalaci\'{o}n de la NSA y una pintura de paisaje generan ambas interacci\'{o}n; ambas alimentan la m\'{a}quina. La tuber\'{\i}a es: produce contenido para el feed $\rightarrow$ s\'{e} descubierto por curadores (que te encontraron en Instagram) $\rightarrow$ exhibe el mismo contenido o una cr\'{\i}tica de \'{e}l en instituciones cuya audiencia tambi\'{e}n fue encontrada en Instagram. \citet{fraser2005critique} identific\'{o} este bucle en la cr\'{\i}tica institucional hace tres d\'{e}cadas; la hegemon\'{\i}a de plataformas es el mismo problema estructural con capitalismo de vigilancia superpuesto.

La falla espec\'{\i}fica de la cr\'{\i}tica-como-contenido no es que opere dentro de la plataforma---todo lo hace---sino que \emph{critica sin cultivar}. Usa el alcance de la plataforma para amplificar un diagn\'{o}stico y se detiene ah\'{\i}. No usa el alcance para preparar un p\'{u}blico para una pr\'{a}ctica que viva en otra pila. La plataforma est\'{a} perfectamente dispuesta a distribuir el diagn\'{o}stico; el diagn\'{o}stico es contenido, y el contenido es lo que la plataforma distribuye. Lo que la plataforma no puede distribuir es un usuariado. Ese es un trabajo que debe hacerse deliberadamente, con una forma que la plataforma pueda llevar pero no reproducir.

\section{Espacios libres de juego}

El arte no crece en feeds. \citet{winnicott1971playing} llam\'{o} al espacio donde la experiencia creativa se hace posible \emph{espacio potencial}---un \'{a}rea transicional de experiencia entre la realidad interior y exterior, fr\'{a}gil, que requiere seguridad, continuidad y libertad de intrusi\'{o}n. \citet{huizinga1938homo} describi\'{o} el ``c\'{\i}rculo m\'{a}gico'' del juego: un espacio acotado con sus propias reglas, separado de la vida ordinaria. Dentro del c\'{\i}rculo m\'{a}gico, las reglas del mundo exterior est\'{a}n suspendidas. Un estudio, cuando funciona, es un c\'{\i}rculo m\'{a}gico. Una sala de ensayo es un c\'{\i}rculo m\'{a}gico. Un cuaderno de bocetos es un c\'{\i}rculo m\'{a}gico.

La plataforma no tiene c\'{\i}rculo m\'{a}gico. Todo acto creativo realizado en Instagram es inmediatamente medido, ordenado, clasificado y devuelto como dato de interacci\'{o}n. El espacio potencial colapsa en m\'{e}tricas en el momento en que la obra es publicada. \citet{fisher2009capitalist} llam\'{o} \emph{realismo capitalista} a la sensaci\'{o}n generalizada de que no hay alternativa: el realismo no est\'{a} en el contenido sino en la incapacidad de imaginar otra cosa. La escuela de arte ense\~{n}a la forma de un c\'{\i}rculo m\'{a}gico; las plataformas no. Este art\'{\i}culo es un argumento a favor de imaginar otra cosa---y un informe sobre lo que esa ``otra cosa'' parece una vez construida.

\section{Anti-entornos}

McLuhan argumentaba que no se puede percibir un entorno desde dentro de \'{e}l~\citep{mcluhan1964understanding, mcluhan1969playboy}. Un entorno es invisible para sus habitantes precisamente porque es total. La \'{u}nica manera de hacer visible un entorno es construir un \emph{anti-entorno}: una contra-estructura cuya diferencia con el entorno dominante revela la forma del entorno---lo que \citet{shklovsky1917art} llam\'{o} \emph{desfamiliarizaci\'{o}n}, la capacidad del arte de hacer extra\~{n}o lo habitual.

Este art\'{\i}culo usa \emph{anti-entorno} en su sentido estructural: no una protesta, no una vitrina, no un rechazo. Una alternativa que funciona. Un instrumento. Una estructura acotada, portadora de reglas, construida junto al entorno dominante y que enruta el trabajo creativo por su propia pila en lugar de por la de la plataforma. La relaci\'{o}n con la plataforma no es ni oposicional ni abstencionista. Un anti-entorno puede usar el alcance de una plataforma sin heredar su l\'{o}gica, de la misma manera que el arte correo us\'{o} el servicio postal sin convertirse en una oficina postal. La radio comunitaria, las redes de arte correo, la cultura \emph{zine}, los circuitos de distribuci\'{o}n Fluxus y el \emph{Computer Lib} de Nelson~\citep{nelson1974computerlib} funcionaron todos como anti-entornos respecto a los medios de difusi\'{o}n de su \'{e}poca; \citet{lanier2018ten} llega a una conclusi\'{o}n afín desde dentro de la industria. Algunos sobrevivieron; muchos fueron absorbidos. La supervivencia depend\'{\i}a menos del rechazo que de si el anti-entorno ten\'{\i}a una pr\'{a}ctica que el entorno dominante no pudiera reproducir.

Un anti-entorno para la computaci\'{o}n creativa en 2026 es una pila: un runtime, un lenguaje, un instrumento, un sustrato de hardware, un protocolo comunitario. Cada capa es un compromiso estructural espec\'{\i}fico que rechaza la l\'{o}gica de plataforma a nivel de la capa, no a nivel del mensaje. La siguiente secci\'{o}n describe los compromisos que \ac{} ha adoptado, cada uno documentado en profundidad por un art\'{\i}culo acompa\~{n}ante en el \emph{platter} de investigaci\'{o}n.

\section{Lo que se ha construido}

Cinco a\~{n}os de operaci\'{o}n han producido un conjunto de compromisos estructurales. Cada uno es un rechazo espec\'{\i}fico de la l\'{o}gica de plataforma a nivel de la pila, y cada uno est\'{a} documentado en profundidad por un art\'{\i}culo hermano en el \emph{platter} de investigaci\'{o}n. Esta secci\'{o}n es la s\'{\i}ntesis: estas no son decisiones separadas. Son el mismo movimiento.

\subsection{URL-primero: la direcci\'{o}n es la interfaz}

Cada pieza en \ac{} es una URL. No hay feed, no hay men\'{u}, no hay selector de archivos, no hay lista de proyectos, no hay superficie algor\'{\i}tmica. El prompt \emph{es} la barra de direcciones: escribes una palabra y est\'{a}s en la pieza~\citep{scudder2026urltradition}. Esta es la tradici\'{o}n HyperCard re-anclada a la web, descendiendo a trav\'{e}s de los enlaces de remix instant\'{a}neo de Glitch y las URL del Web Editor de p5.js. Una plataforma tiene una p\'{a}gina de inicio a la que llegas; \ac{} tiene cien direcciones hacia las que escribes. La URL no es una funcionalidad. Es la propiedad-medio de todo el sistema.

\subsection{Pieza-no-publicaci\'{o}n: la unidad acotada}

La unidad de producci\'{o}n cultural en \ac{} es una \emph{pieza}: un programa interactivo autocontenido~\citep{scudder2026pieces}, direccionable por nombre, persistente en el tiempo, no ordenado algor\'{\i}tmicamente. Una pieza es a una publicaci\'{o}n lo que una canci\'{o}n es a un scroll. No descubres una pieza porque un algoritmo la aflor\'{o}; llegas a una pieza porque alguien te la nombr\'{o}, o porque escribiste su direcci\'{o}n, o porque la encontraste en la lista p\'{u}blicamente buscable. El \emph{espacio potencial} de Winnicott~\citep{winnicott1971playing} se restaura a nivel de la unidad: cada pieza est\'{a} acotada, reglada y no interrumpida por m\'{e}tricas de interacci\'{o}n.

\subsection{Instrumento-no-perfil: la interfaz se toca}

\np{}~\citep{scudder2026notepat} es un sintetizador de teclado crom\'{a}tico que corre en el navegador y en el hierro desnudo. Es un instrumento que tocas. No es un perfil que curas, no es un feed que scroleas, no es una marca que mantienes. La diferencia entre instrumento y perfil es la diferencia entre uso y exhibici\'{o}n. Un perfil est\'{a} optimizado para ser visto; un instrumento est\'{a} optimizado para ser usado. \ac{} trata la interfaz creativa como un instrumento en cada nivel de la pila---desde el teclado de \np{} al REPL de KidLisp al propio prompt---y esta elecci\'{o}n degrada la legibilidad-para-otros como objetivo de dise\~{n}o a favor de la inmediatez-para-uno-mismo.

\subsection{KidLisp: un lenguaje de la era de la composici\'{o}n}

KidLisp~\citep{scudder2026kidlisp} es un dialecto m\'{\i}nimo de Lisp para arte generativo. Es un lenguaje de la era de la composici\'{o}n en el sentido de Attali~\citep{attali1985noise}: la distinci\'{o}n entre int\'{e}rprete y audiencia se disuelve dentro de \'{e}l. El mismo c\'{o}digo que escribes para pintar un cuadro es el c\'{o}digo que otros leen para entender el cuadro. El mismo REPL que usas para bocetar es el REPL que otros usan para bifurcar. Esta no es una propiedad general de Lisp---la mayor\'{\i}a de los Lisps son de la era del arquitecto, construidos para especificar sistemas---sino una propiedad de KidLisp espec\'{\i}ficamente, que cambia potencia expresiva por inmediatez, forkabilidad y legibilidad a resoluci\'{o}n de 8 bits.

\subsection{Nodo soberano: el hardware es tuyo}

\ac{} se env\'{\i}a como un sistema operativo de hierro desnudo que arranca hardware excedente directamente en el runtime creativo~\citep{scudder2026os, scudder2026identity}. No hay cuenta en la nube, no hay bloqueo de proveedor, no hay temporizador de expiraci\'{o}n de actualizaciones, no hay telemetr\'{\i}a de comportamiento. El port\'{a}til es un nodo soberano en la red: ejecuta su propio servidor, guarda sus propios datos, federa con otros nodos en sus propios t\'{e}rminos. Esto es la anti-Chromebook~\citep{scudder2026openschools}: el mismo hardware de 50 d\'{o}lares, flasheado con software que el usuario posee, sobre una pila de protocolos que puede auditar. Con aproximadamente 240 millones de m\'{a}quinas Windows 10 en fin de vida \'{u}til yendo al vertedero en el ciclo 2025--2026, el compromiso del nodo soberano no es un gesto pol\'{\i}tico de nicho; es un programa material para qu\'{e} hacer con el hardware.

\subsection{Pals: afinidad, no seguidores}

La comunidad de \ac{} est\'{a} organizada seg\'{u}n el modelo Pals, heredado de las giras de Goodiepal y documentado en profundidad en el art\'{\i}culo sobre Goodiepal~\citep{scudder2026goodiepal}. Un Pal no es un seguidor. Un Pal es alguien que hace contigo, viaja contigo, comparte el instrumento contigo. La relaci\'{o}n es rec\'{\i}proca, basada en la afinidad, y cara a cara donde es posible; no es un canal de difusi\'{o}n unidireccional. Las plataformas optimizan para el flujo parasocial uno-a-muchos porque eso es lo que la publicidad monetiza. Los Pals optimizan para la profundidad uno-a-uno porque eso es lo que la pr\'{a}ctica recompensa. Los n\'{u}meros son menores. Los v\'{\i}nculos son m\'{a}s densos. El modelo escala por bifurcaci\'{o}n, no por alcance. ``Making kin with the machines'' de Lewis~\citep{lewis2018making}, \emph{Dise\~{n}os para el pluriverso} de \citet{escobar2018designs} y la justicia de dise\~{n}o de \citet{costanzachock2020design} articulan el mismo compromiso en registros diferentes: relacional m\'{a}s que transaccional, cuidado m\'{a}s que extracci\'{o}n.

\subsection{El linaje de 94 proyectos}

Nada de lo anterior es nuevo. La arqueolog\'{\i}a del repositorio de \ac{} traza la genealog\'{\i}a t\'{e}cnica de la plataforma a trav\'{e}s de 94 proyectos predecesores que abarcan 2007--2020~\citep{scudder2026archaeology}---herramientas de dibujo, herramientas de performance, infraestructura de desarrollo, juegos, instrumentos. Los compromisos de arriba no llegaron como filosof\'{\i}a de dise\~{n}o. Se acumularon como resoluci\'{o}n de problemas que seguían reapareciendo a lo largo de 94 intentos. El anti-entorno no es una idea que el autor tuvo. Es la forma que cay\'{o} de practicar la misma pr\'{a}ctica durante dos d\'{e}cadas dentro de entornos medi\'{a}ticos que eran, cada uno a su vez, totales.

\section{El motor y el usuariado}

La objeci\'{o}n al argumento hasta ahora se escribe sola: si los anti-entornos se construyen fuera de la plataforma, ¿qu\'{e} hacer con la cuenta de TikTok @whistlegraph, que alcanz\'{o} 2,6 millones de seguidores durante la pandemia de COVID-19~\citep{scudder2026whistlegraph} y todav\'{\i}a mueve aproximadamente cuarenta y ocho mil visualizaciones de v\'{\i}deo al d\'{\i}a m\'{a}s de dos a\~{n}os despu\'{e}s de que el tr\'{\i}o se separara en noviembre de 2023?\footnote{Instant\'{a}nea de TikTok Analytics de @whistlegraph, del 24 de abril de 2025 al 23 de abril de 2026---el segundo a\~{n}o completo tras la separaci\'{o}n del tr\'{\i}o, durante el cual no se produjeron nuevos whistlegraphs por parte de los creadores originales: $\sim$17,69 millones de visualizaciones de v\'{\i}deo (mediana 43\,047/d\'{\i}a), $\sim$465\,000 likes, $\sim$16\,000 compartidos, $\sim$191\,000 visitas de perfil. Los picos diarios (208\,000 el 27 de septiembre de 2025; 198\,000 el 16 de febrero de 2026) reflejan circulaci\'{o}n org\'{a}nica continuada del archivo original.} ¿No es eso prueba de que el exterior s\'{o}lo existe con permiso del interior?

Es prueba de algo m\'{a}s preciso. Las plataformas son motores de distribuci\'{o}n. Mueven atenci\'{o}n a una velocidad y escala que ninguna tuber\'{\i}a institucional se acerca, y de forma suficientemente fiable como para que un artista pueda tratarlas como infraestructura. La pregunta que la literatura est\'{a}ndar de cr\'{\i}tica de plataformas no se hace---porque est\'{a} comprometida con un marco oposicional---es qu\'{e} est\'{a} llevando el motor. La cuenta @whistlegraph llevaba una gram\'{a}tica formal espec\'{\i}fica: notaci\'{o}n gr\'{a}fica (marcas que corresponden a s\'{\i}labas cantadas), ret\'{o}rica de interfaz (performance como conferencia como instrumento) y una l\'{o}gica de reproducibilidad (la partitura que ense\~{n}a c\'{o}mo tocarla~\citep{scudder2026whistlegraph}). Esa gram\'{a}tica se desarroll\'{o} pre-plataforma a trav\'{e}s de Radical Digital Painting~\citep{scudder2017manifesto, scudder2018rdp35c3}, 65 conferencias-performance incluyendo una gira europea en 2018 con Goodiepal~\citep{scudder2026goodiepal}, y una d\'{e}cada de pr\'{a}ctica dibujo-como-canto en la que la plataforma no tuvo parte en escribir. TikTok no hizo el whistlegraph; distribuy\'{o} uno.

Lo que la distribuci\'{o}n prepar\'{o} es un \emph{usuariado}: una poblaci\'{o}n de aproximadamente 2,6 millones de personas que, al ver y reproducir whistlegraphs, aprendieron a reconocer un tipo espec\'{\i}fico de movimiento formal---marcas como instrucciones, interfaces como instrumentos, performance como pedagog\'{\i}a. El usuariado no es el conteo de seguidores. Es la audiencia preparada disponible aguas abajo, al ritmo en que pueda convertirse en pr\'{a}ctica sobre un sistema que lleva la misma gram\'{a}tica m\'{a}s profundamente de lo que la plataforma puede albergar. \ac{} es ese sistema~\citep{scudder2026ac}. El usuariado es su poblaci\'{o}n-semilla. La auditor\'{\i}a de red~\citep{scudder2026audit} documenta la forma actual del usuariado convertido---2.800+ handles registrados, 16.000+ programas KidLisp, 93.000+ sesiones de arranque---una peque\~{n}a fracci\'{o}n del p\'{u}blico preparado y m\'{a}s de lo que cualquier herramienta comparable de computaci\'{o}n creativa ha producido sin un progenitor institucional.

La paradoja que el borrador anterior de este art\'{\i}culo no pod\'{\i}a resolver es esta: el anti-entorno m\'{a}s grande en el \emph{platter} de \ac{} fue construido \emph{a trav\'{e}s} del complejo. La resoluci\'{o}n es que el complejo nunca fue el adversario. Era el motor. El error que comete la literatura de cr\'{\i}tica-como-contenido es leer cada uso de la plataforma como capitulaci\'{o}n. El movimiento real del artista---el disponible para cualquiera con una forma que el algoritmo de la plataforma amplificar\'{a}---es usar el motor deliberadamente, entregando una carga que la plataforma no puede inspeccionar, hacia un usuariado que la plataforma no puede albergar. La escala se alquila al motor. La pr\'{a}ctica la posee el usuariado.

\section{Financiar el motor y el c\'{\i}rculo}

Los anti-entornos son dif\'{\i}ciles de financiar por una raz\'{o}n estructural: la econom\'{\i}a de plataformas pone precio a la atenci\'{o}n, y el anti-entorno est\'{a} construido para rechazar la mensurabilidad de la atenci\'{o}n. Una pieza no tiene tasa de interacci\'{o}n. Un Pal no tiene coste por mil. Un instrumento no tiene curva de retenci\'{o}n. El art\'{\i}culo sobre sostenibilidad en esta serie~\citep{scudder2026sustainability} revisa c\'{o}mo se han financiado 28 herramientas comparables de computaci\'{o}n creativa---subsidio acad\'{e}mico, mecenazgo corporativo, sacrificio personal, crowdfunding comunitario, golpe de suerte o ning\'{u}n modelo en absoluto---y encuentra que las herramientas m\'{a}s usadas no son las mejor financiadas. La brecha mediana entre la creaci\'{o}n de una herramienta y su primera financiaci\'{o}n sostenible es de ocho a\~{n}os. La mayor\'{\i}a de las herramientas cierran antes de cruzarla.

El marco motor-y-usuariado no resuelve el problema de financiaci\'{o}n, pero lo clarifica. El motor es gratuito para usar y gratuito para dejar. El usuariado es el activo que el motor produce. La pr\'{a}ctica es donde se gasta el activo. La pregunta de financiaci\'{o}n no es si monetizar el motor---el motor monetiza a los artistas en sus propios t\'{e}rminos, vendiendo su atenci\'{o}n a los anunciantes---sino c\'{o}mo convertir el usuariado en una econom\'{\i}a de pr\'{a}ctica autosostenible antes de que la pr\'{o}xima actualizaci\'{o}n del algoritmo del motor rote la atenci\'{o}n hacia otra parte. Las \emph{herramientas para la convivialidad} de Illich~\citep{illich1973tools} nombran el destino: herramientas construidas a escala humana, a mano, para el uso en lugar del consumo. La econom\'{\i}a de tales herramientas todav\'{\i}a se est\'{a} escribiendo. Este art\'{\i}culo es una entrada en su documentaci\'{o}n.

\section{Coda}

Los anti-entornos son posibles. No son raros; simplemente no son legibles para la plataforma, porque la plataforma no fue construida para leer lo que producen. Este art\'{\i}culo ha descrito la forma de uno, desde dentro, con los compromisos estructurales espec\'{\i}ficos sobre los que descansa puestos por escrito para que otros puedan reconocerlos, bifurcarlos, adaptarlos o rechazarlos:

\begin{itemize}
  \item URL-primero, no feed-primero.
  \item Pieza-no-publicaci\'{o}n: unidades acotadas, no flujos de scroll.
  \item Instrumento-no-perfil: interfaces para tocar, no para curar.
  \item Lenguaje de la era de la composici\'{o}n: c\'{o}digo le\'{\i}do como se escribe.
  \item Nodo soberano: hardware y datos que posees.
  \item Afinidad, no seguidores: escala por bifurcaci\'{o}n, no por alcance.
  \item Motor, no adversario: usa el alcance de la plataforma para preparar un usuariado cuya pr\'{a}ctica vive en otra parte.
\end{itemize}

Estas son las reglas de un c\'{\i}rculo. Otras personas pueden dibujar los suyos. El algoritmo no puede leer ninguna de ellas.

\vspace{0.5em}
\noindent\textbf{ORCID:} \href{https://orcid.org/0009-0007-4460-4913}{0009-0007-4460-4913}

\vspace{0.5em}
\noindent\textit{Traducido del ingl\'{e}s. Versi\'{o}n original disponible en \url{https://papers.aesthetic.computer}}

\bibliographystyle{plainnat}
\bibliography{references}

\end{document}
